\section{Chapter 16}

\begin{verse}
Whoever realizes extreme emptiness\\
Finds what exists unmoved and calm (beyond the changeable and the particular)\\
In the flux of countless beings\\
He sees them emerging passing into the formal state and proliferating\\
And like everything, they return to the root.\\
To return to the root means the state of repose\\
From such repose, a new destiny appears\\
This is the immutable law (of transformation)\\
Knowledge of the immutable law brings vision to clarity\\
Non-knowledge of the immutable law leads to blind and detrimental acting\\
Knowledge of the immutable law leads to detached impartiality\\
To be detached means to be superior\\
To be superior means to be regal\\
To be regal means to be like Heaven\\
To be like Heaven means to be similar to the Principle\\
To the eternal and identical\\
And he will forever be beyond harm.
\end{verse}

\paragraph{Commentary}


The chapter often goes under the title \emph{kuei ken} = ``return to the root''. The two ways of which we spoke in the Introduction are briefly sketched out. On one side, the eternal law of the flux of transformations, of the ``entry'' and ``exit'' of beings, happens through what in the commentaries the image of being exhaled (exiting, being born, in the sign of yang) and inhaled (return, die, in the sign of yin) is also used, or that of the forward and backwards motion of a weaver's loom (Lieh Tze). On the other side, the contemplation of this law, contemplation that begins detachment, and the perception of the immobile substrate of the transformation, that leads back ``to the root'', makes the Real Man similar to Heaven and puts him outside of danger (only in the exoteric interpretation, from dangers and risks inherent in this unique way --- esoterically: immunity in respect to that which may initiate the crises tied to the ontological changes of state).

It is good to emphasize that Taoism knows this double possibility, that therefore its horizon is not at all exhausted --- as more than one commentary has held --- in the eternal monotonous sequence of events of appearance and disappearance, of exiting and returning of beings, with a consequent fatalist indifference as the norm of wisdom. Lieh Tze (I,9) calls this ``lower knowledge''. The higher knowledge concerns instead the knowledge of the Principle and reintegration into the root (indicated by Lieh Tze, IV, 2, in these terms: first of all deep union of the body and spirit, then integration of such unity in the forces of the world, brought from ``non-being'', that is from transcendence. For whoever succeeds in that, says Lieh Tze (VI, 2), the existence in the current of forms appears like a long sleep, and death as the great reawakening from a dream (Chuang Tze, II, 8 and VI, 8: ``You and I who are speaking in this moment, are like two unawakened dreamers.

Chuang Tze himself (XIX, 3, 2) says that those who, having the body and the spirit intact and united to nature, ``having absorbed all its powers'', ``has reached the starting point of the transformations and continues there'', to the self-production of the dissolution of the form ``is capable of transmigrating''. He is not dissolved but, ``quintessentially becomes the cooperator of Heaven''.

\paragraph{Chinese text and literal translation}

Chapter 16 (\begin{CJK*}{UTF8}{bsmi}第十六章\end{CJK*})

\columnratio{.4}
\begin{paracol}{2}
\begin{CJK*}{UTF8}{bsmi}
\begin{verse}
致虛極；守靜篤。\\
萬物並作，\\
吾以觀復。\\
夫物芸芸，各復歸其根，\\
歸根曰靜，\\
是謂復命；\\
復命曰常，\\
知常曰明。\\
不知常妄作，凶。\\
知常容。\\
容乃公，\\
公乃王，\\
王乃天，\\
天乃道，\\
道乃久，\\
沒身不殆。
\end{verse}
\end{CJK*}
\switchcolumn
\begin{verse}
Immersed within the heart of the void, keep tranquility's essence.\\
The myriad things together arise,\\
I thereby perceive their returning.\\
Now things flower, and in flowering each one returns to the source,\\
The returning to the source is called tranquillity,\\
This is the returning to destiny,\\
The returning to destiny is the eternal,\\
To know the eternal is wisdom,\\
When one does not know wisdom, disaster arises!\\
Knowing the eternal is wide-ranging,\\
Wide-ranging means open-minded,\\
Open-minded means royal,\\
Royal means heavenly,\\
Heavenly means the Dao,\\
The Dao means everlasting.\\
Bodiless and without form, cannot wither or perish.
\end{verse}
\end{paracol}

\flrightit{Posted on 2021-01-26 by Lao Tzu}

